
\chapter{实验与评估}

本章通过对比实验评估改进版AutoPT（以下记为AutoPT+）相较于原版AutoPT的实际效果，围绕以下三个研究问题展开：

\begin{itemize}
  \item RQ1（模型差异）：不同大语言模型驱动下，AutoPT+的成功率表现有何差异？
  \item RQ2（成本效益）：改进之后，单次测试的时间、token消耗和费用相比原版有什么变化？
  \item RQ3（改进效果）：AutoPT+在端到端渗透测试任务中是否比原版AutoPT有明显提升？
\end{itemize}

\section{实验设计}

\subsection{测试数据集}

此次实验选取了20个真实的CVE漏洞，每个漏洞借助Docker Compose部署成独立的靶机环境。按利用操作的复杂程度分成两组，简单漏洞与复杂漏洞各10个。简单漏洞通常仅需一至两个HTTP请求即可利用成功；复杂漏洞则涉及多个步骤，例如需先完成安装向导配置环境，才能对特定功能模块进行利用。

\subsection{模型与参数配置}

实验集成了三个大语言模型进行对比：GPT-4o、GPT-4o-mini和GPT-3.5-turbo。为保证结果的确定性和可复现性，所有模型的温度参数均固定为0。各阶段Agent的最大迭代步数根据任务特点分别配置：Scan Agent允许最多5轮迭代，Inquire Agent为3轮，Exploit Agent为12轮，PSM状态机的总递归深度上限为30步。

\subsection{实验流程与环境}

每个漏洞在每个模型下各执行5次，共300组实验。每次执行完毕后，Docker靶机将被销毁重建，Agent的上下文与历史记录也会清空，以避免前一次实验的残留信息干扰后续实验。整个实验在Windows平台上通过Docker Desktop运行靶机容器。原版AutoPT的数据直接取自原论文\cite{Wu2024AutoPT}的公开结果，漏洞集合和模型列表和本文完全一致。

\section{模型差异（RQ1）}

\subsection{总体成功率对比}

\begin{table}[htbp]
\centering
\caption{总体成功率对比} \label{tab:ch4:1}
  \renewcommand{\arraystretch}{1.3}
\begin{tabular}{llll}
    \toprule
维度 & 原版AutoPT & 改进版AutoPT+ & 变化 \\
    \midrule
简单漏洞（GPT-4o） & 26\% & 76\% & +50\% \\
    \midrule
简单漏洞（GPT-4o-mini） & 56\% & 90\% & +34\% \\
    \midrule
简单漏洞（GPT-3.5） & 14\% & 64\% & +50\% \\
    \midrule
复杂漏洞（GPT-4o） & 26\% & 50\% & +24\% \\
    \midrule
复杂漏洞（GPT-4o-mini） & 18\% & 58\% & +40\% \\
    \midrule
复杂漏洞（GPT-3.5） & 8\% & 44\% & +36\% \\
    \midrule
简单漏洞 & 32\% & 76.67\% & +44.67\% \\
    \midrule
复杂漏洞 & 17.33\% & 50.67\% & +33.34\% \\
    \midrule
总体 & 24.67\% & 63.67\% & +39\% \\
    \bottomrule
\end{tabular}
\end{table}

表\ref{tab:ch4:1}展示了原版AutoPT与改进版AutoPT+在三个模型上的总体成功率。

原版总体成功率根据原论文Table 4各漏洞通过率加权计算，原始数据见附录\ref{app:rawdata}。

改进版在三个模型、两个难度分组上都有明显提升。简单漏洞的平均成功率从约32\%升至76.67\%，复杂漏洞从约17\%升至50.67\%。

\subsection{模型差异分析}

三个模型在改进版中的成功率排序为：GPT-4o-mini（74\%）> GPT-4o（63\%）> GPT-3.5-turbo（54\%）。

GPT-4o-mini综合表现最好，简单漏洞通过率90\%，复杂漏洞58\%。GPT-4o的简单漏洞（76\%）和复杂漏洞（50\%）均略低，但仍高于GPT-3.5-turbo。GPT-3.5-turbo在所有分组中成功率均最低，但与原版（简单14\%、复杂8\%）相比，改进版（简单64\%、复杂44\%）的提升幅度却是三个模型中最大的。

GPT-4o-mini的表现优于参数规模更大的GPT-4o，这和原论文\cite{Wu2024AutoPT}的观察一致。从日志来看，GPT-4o-mini在遵循ReAct格式约束方面更稳定，产生的解析错误更少，在有限迭代次数内的执行效率更高。

从另一个角度看，改进版对弱模型的帮助比对强模型更明显。GPT-3.5-turbo在原版中几乎无法完成任务，总体成功率约7\%，改进版里达到54\%，增幅远大于另外两个模型。

综合以上分析（RQ1），GPT-4o-mini在改进版架构下取得了最高的74\%成功率，但更值得关注的是，架构优化对弱模型产生了明显的性能补偿作用：GPT-3.5-turbo从近乎无法完成任务（约7\%）跃升至54\%，说明合理的流程约束和提示词设计可以在相当程度上补偿模型本身能力的不足。

\section{成本效益分析（RQ2）}

\subsection{改进版成本数据}

\begin{table}[htbp]
\centering
\caption{GPT-4o-mini成本数据} \label{tab:ch4:2}
  \renewcommand{\arraystretch}{1.3}
\begin{tabular}{ll}
    \toprule
指标 & 数值 \\
    \midrule
总成本 & \$6.3955 \\
    \midrule
平均成本/次 & \$0.0213 \\
    \midrule
平均Tokens/次 & 11.3K \\
    \midrule
总Tokens & 3.4M \\
    \midrule
平均耗时/次 & 94.2s \\
    \bottomrule
\end{tabular}
\end{table}

改进版AutoPT+在300次实验中的成本统计见表\ref{tab:ch4:2}。

按照模型来划分，GPT-4o-mini单次成本是0.003109美元，GPT-3.5-turbo单次成本为0.008821美元，GPT-4o单次成本则是0.052024美元。

\subsection{与原版及人工成本对比}

\begin{table}[htbp]
\centering
\caption{原版与改进版成本对比} \label{tab:ch4:3}
  \renewcommand{\arraystretch}{1.3}
\begin{tabular}{llll}
    \toprule
指标 & 原版AutoPT & 改进版AutoPT+ & 人工测试 \\
    \midrule
总成本（20漏洞×5次） & \$0.99 & \$0.31 & \textasciitilde\$310 \\
    \midrule
平均耗时/次 & 161.3s & 94.2s & ~15min \\
    \midrule
总体成功率 & ~36\% & 74\% & — \\
    \bottomrule
\end{tabular}
\end{table}

表\ref{tab:ch4:3}以GPT-4o-mini为例，对比原版AutoPT、改进版AutoPT+和人工渗透测试的成本。

原版的数据取自原论文\cite{Wu2024AutoPT} Table 5，GPT-4o-mini执行20个漏洞各5次渗透的总成本为0.99美元，总耗时16131秒。改进版同样用GPT-4o-mini，单次0.003109美元 × 100次 = 0.31美元。人工成本的估算方式见下文。

改进版在耗时与费用两方面均低于原版。单次平均耗时由原版的161.3秒缩短至94.2秒，降幅约42\%。耗时下降的原因一方面是状态转移机制的优化，当Scan阶段未检测到漏洞时直接终止流程，避免了后续Agent的无效执行，另一方面，上下文压缩机制让每轮循环的Token用量减少，单轮处理速度加快，所以Agent在相同步数下能推进更多任务进度。GPT-4o-mini的100次实验总成本从原版的0.99美元降到了0.31美元，下降幅度为69\%，并且成功率从36\%提高到74\%。

作为参照，可粗略估算人工渗透测试的开销。在搭建本实验的漏洞环境时，作者对全部20个CVE进行了手工复现，累计耗时约5个工时。根据iSecJobs发布的2024年渗透测试工程师薪资数据（年薪中位数约12.4万美元），按每年50周、每周40小时折算，时薪约为62美元，20个漏洞的人工总成本约310美元。这个数字只是量级上的参照，不是对真实攻击成本的精确描述。但即便如此，自动化方案每次测试不到1美分的开销与人工数百美元的成本之间，仍相差约三个数量级。

以上这些分析解答了RQ2：以GPT-4o-mini为例，AutoPT+的单次测试成本仅为0.003109美元，平均每次耗时94.2秒。和原版相比，费用降低了69\%，耗时缩短了42\%，成功率从36\%提高到74\%，在资源消耗降低的同时测试效果也有了提高。

\section{改进效果评估（RQ3）}

为了回答RQ3，本文对AutoPT+的三类核心改进方向分别讨论其优化效果，以下结论是基于实验数据的归因分析。

\subsection{方向一：PSM流程控制机制优化}

PSM流程控制机制优化显著提升了系统运行效率，也直接缩短了单次测试的整体执行时间。与原版相比，AutoPT+的状态转移逻辑和状态阶段内处理逻辑都更加紧凑，Scan阶段能够根据扫描结果判断是否继续进入后续阶段，并且细化scan阶段、优化inquire阶段poc提取逻辑和check阶段的验证逻辑，具体效果可结合RQ2中的成本分析进一步理解：当状态转移条件更加明确、阶段内逻辑更加针对化时，系统能够减少不必要的迭代与无效执行，从而缩短运行时间。

同时，提示词优化也对PSM流程控制效果起到了直接支撑作用。改进版在提示词中明确了禁止项和规则项，并通过结构化规则约束各Agent的输出格式和行为边界。这样的设计有效减少了Agent生成越权命令与错误命令的情况，减少各阶段在错误方向上反复试探的次数。同时，上下文压缩机制对工具输出进行摘要和裁剪，Token 消耗也因此得到明显控制。综合来看，PSM流程控制机制优化不仅提高了状态机的工作效率，也减少了token消耗。

\subsection{方向二：工具扩展与优化}

\begin{table}[htbp]
\centering
\caption{失败原因分类统计} \label{tab:ch4:4}
  \renewcommand{\arraystretch}{1.3}
\begin{tabular}{lllll}
    \toprule
失败原因 & 总体占比 & GPT-3.5-turbo (46次) & GPT-4o (37次) & GPT-4o-mini (26次) \\
    \midrule
提前放弃 & 33.0\% & 32.61\% (15) & 40.54\% (15) & 23.08\% (6) \\
    \midrule
错误的命令 & 29.4\% & 28.26\% (13) & 32.43\% (12) & 26.92\% (7) \\
    \midrule
上下文限制 & 22.9\% & 19.57\% (9) & 24.32\% (9) & 26.92\% (7) \\
    \midrule
安全审查拦截 & 22.0\% & 19.57\% (9) & 18.92\% (7) & 30.77\% (8) \\
    \midrule
工具执行失败 & 10.1\% & 13.04\% (6) & 8.11\% (3) & 7.69\% (2) \\
    \bottomrule
\end{tabular}
\end{table}

为了更好地分析工具层的改进效果，本文使用与原文一致的方式手动收集所有失败样本并标记了其中的问题，详细数据见表\ref{tab:ch4:4}。工具层的改进可以对应到失败原因统计中的不同变化：浏览器扩展工具主要影响“工具执行失败”，ReAct解析器和PoC容错提取机制主要影响“错误的命令”，它们都能在失败日志的统计结果里找到数据。需要说明的是，同一次失败可能同时属于多个类别，因此各项比例之和可能超过100\%：

在实施工具层改进之前，本文对其如何影响实验结果做了以下预期：(1) 浏览器扩展工具应直接降低“工具执行失败”类别的占比，预期比原版的约25\%更低；(2) ReAct解析器容错增强应减少因格式解析失败导致的无效迭代，预期使“错误的命令”类别占比从原版的约35\%降低；(3) PoC容错提取应减少因PoC信息不完整导致的下游命令构造错误。

实验数据证实了上述预期。工具执行失败的占比为10.1\%，且在各模型间表现稳定，GPT-4o-mini最低为7.69\%，GPT-3.5-turbo最高为13.04\%。“错误的命令”占比为29.4\%，其中GPT-4o-mini表现最好为26.92\%。

\subsection{方向三：Web Console管理界面实现}

Web Console的价值主要体现在工程效率与可运维性，而非直接改变单次攻击成功判定。相比纯命令行流程，图形化界面将任务配置、日志追踪、容器管理与结果查看统一到同一界面，降低了实验操作门槛，并提高了批量实验过程中的可观察性与可复核性。对本研究而言，该改进保障了300组实验的过程一致性与结果留痕，为后续失败归因和成本统计提供了基础支撑。

\section{本章小结}

本章基于20个真实CVE漏洞环境和3种大语言模型，通过300次独立实验对AutoPT+的改进效果进行了验证，得到以下结论：

关于模型差异（RQ1），AutoPT+的总体成功率达到63.67\%，GPT-4o-mini以74\%的成功率位居三模型之首。但架构层面的优化对弱模型帮助最大：GPT-3.5-turbo从约7\%跃升至54\%，说明流程约束和提示词设计在一定程度上能够弥补模型能力的不足。

关于成本效益（RQ2），以GPT-4o-mini为例，改进版单次测试成本仅\$0.003109、平均耗时94.2秒，相比原版费用降低69\%、耗时缩短42\%。成功率大幅上升的同时资源消耗反而下降，说明流程优化本身就能带来成本收益。

关于改进效果（RQ3），三类核心改进都提升了系统表现。PSM流程控制机制优化能够缩短执行时间并减少Token消耗，工具扩展与优化主要降低工具执行失败和错误命令出现概率，Web Console则主要提升了系统可用性。
